\silentchapter{6 Forma oppstår mellom agentane}{forma oppstår mellom agentane}
\prop{6}{}{Forma oppstår mellom agentane.}
\prop{6.1}{t}{Av 2.12 og 5.1: sidan meir enn eitt seleksjonstrykk alltid verkar, og sidan ulike agentar responderer på ulike trykk, er kvar realisert form eit poly-agentisk kompromiss.\footnote{Odling-Smee, Laland \& Feldman (2003)\footnote{Eit komplekst distribuerte system si åtferd kan ikkje deduserast frå kjeldekoden til éin av nodane; forma spring ut av kompromiss i nettverket.}}}
\prop{6.11}{t}{Materialaffordansen, reiskapen, marknaden og kulturen utgjer uavhengige gradientar\footnote{Brukaropplevinga emergerer presis der forretningslogikk, databaseskjema og nettverkslatens kryssar kvarandre som gradientar.}.}
\prop{6.111}{t}{Ingen einskild agent dikterer den resulterande morfologien.}
\prop{6.1111}{t}{Ho emergerer i skjeringspunktet mellom dei.}
\prop{6.11111}{t}{Kvart hierarkisk nivå er ein kompetent problemløysar innanfor sin eigen lyskjegle.}
\prop{6.11112}{t}{Dei lågare komponentane navigerer etter lokale gradientar.}
\prop{6.111121}{t}{Dei bidreg samstundes til mål dei sjølve manglar kapasitet til å representere\footnote{Mikrokomponentar manglar representasjon av heilskapen. Dei leverer data over port 80, men bidreg uvitande til det makroskopiske formrommet vi kallar 'systemet'.}.}
\prop{6.11113}{t}{Reduksjonisme og holisme er komplementære skildringsnivå.}
\prop{6.111131}{t}{Begge er sanne, men ingen er tilstrekkeleg åleine.}
\prop{6.12}{t}{Sjølv når éin agent har nominell kontroll, avgjer landskapet kva reglar som faktisk fyrer.}
\prop{6.121}{t}{Agenten vel frå naboskapen; landskapet filtrerer naboskapen\footnote{Eit strengt typa grensesnitt filtrerer reglane frå grammatikken slik at agenten (brukaren) ikkje feilar; landskapet disiplinerer naboskapet.}.}
\prop{6.1211}{t}{Kontrollen er alltid delt.}
\prop{6.12111}{o}{Kommunikasjonen mellom agentar vert avgrensa av ein endeleg bandbreidde.}
\prop{6.121111}{o}{Ved låg bandbreidde er navigasjonen grov og stokastisk.}
\prop{6.121112}{o}{Ved høg bandbreidde kan forma artikulerast gjennom eksplisitt spesifikasjon\footnote{Med uendeleg bandbreidde kan komponentar dele fulle representasjonar av sine eigne tilstandar, noko som aukar presisjonen proporsjonalt.}.}
\prop{6.12112}{t}{Det finst eit byttetilhøve mellom bandbreidde og uavhengigheit\footnote{Tenesteorientert arkitektur (SOA) demonstrerer at å frikople agentar reduserer koordineringsevna, men sikrar autonom lokal overleving.}.}
\prop{6.121121}{t}{Tett kopla agentar konvergerer mot eit felles representasjonsformat som avgrensar kva regionar dei kan utforske.}
\prop{6.121122}{t}{Lauskopla agentar opprettheld kognitiv autonomi, men koordinerer dårlegare.}
\prop{6.12113}{t}{Det optimale koplingsregimet avheng av kva fase systemet er i.}
\prop{6.121131}{t}{Utforsking krev lauskopla agentar; utnytting krev tett kopla agentar.}
\prop{6.13}{o}{Bandbreidda har ein temporal dimensjon i form av latens\footnote{Asynkronisering løyser bandbreiddeproblemet nettopp ved å innføre tidsforsinka modellar av avhengigheiter.}.}
\prop{6.131}{t}{Ved høg latens må kvar agent operere på ein intern modell av dei andre agentane sin framtidige tilstand.}
\prop{6.1311}{t}{Koordineringa er då berre so presis som modellens adekvatheit.}
\prop{6.13111}{t}{Latensen er ikkje berre ein teknisk eigenskap ved kanalen, men ein epistemisk eigenskap ved agentens situasjon.}
\prop{6.2}{o}{Når koplinga mellom agentar er robust, emergerer eit overordna system.}
\prop{6.21}{t}{Det overordna systemet har større lyskjegle og meir samansette mål.}
\prop{6.211}{t}{Når samanbindinga er tilstrekkeleg tett til at heilskapen realiserer ein eigen tilbakekoplingssløyfe, vert heilskapen sjølv ein agent etter definisjonen i 5.11.\footnote{Kuhn (1962)\footnote{Når fleire skytjenester bindast saman av ein meldingsbuss, vert heilskapen brått sjølv ein meta-agent med sin eigen tilbakekoplingssløyfe.}}}
\prop{6.2111}{o}{Eit multi-agent-system kan oppnå metastabilitet på tvers av fleire skalaer utan at nokon einskild skala er fullt optimalisert.}
\prop{6.21111}{o}{Denne stabiliteten let tradisjonar persistere gjennom fluktuasjonar som ville ha utsletta delane deira om dei opererte i isolasjon.}
\prop{6.3}{t}{Stilar og tradisjonar er deskriptive mellomkonstruksjonar utan kausal kraft.}
\prop{6.31}{t}{Dei skildrar dei punkta i formrommet der vandringa mellombels har stansa.}
\prop{6.311}{t}{Å formgje «i ein stil» er ein logisk feil.}
\prop{6.3111}{t}{Det er å imitere eit statistisk symptom utan å forstå dei seleksjonstrykka som genererte klynga i utgangspunktet.}
\prop{6.31111}{t}{Den som imiterer ein stil, imiterer ein effekt og ikkje ei årsak\footnote{Å kopiere eit populært brukargrensesnitt frå ein annan bransje utan omsyn til dei bakenforliggjande trykka, produserer inkompetent form i den nye klassen.}.}
\prop{6.311111}{t}{Resultatet er ei form som liknar på tidlegare former utan å svare på dei same trykka.}
\prop{6.4}{t}{Nyhetsraten i formrommet er ein funksjon av det tilstøytande moglege.}
\prop{6.41}{o}{Kvar realisert form opnar nye regionar av naboskapen som ikkje fanst før\footnote{Kvart vellukka programvare-API dannar ein plattform, og kvar plattform opnar eit massivt rom av nye, hittil uante moglege tenester.}.}
\prop{6.411}{t}{Det moglege veks med det realiserte.\footnote{Kauffman (1993)}}
\prop{6.4111}{t}{Formhistoria er ikkje ein uttapping av eit gjeve repertoar.}
\prop{6.41111}{t}{Ho er ein ekspansjon av repertoaret gjennom si eiga rørsle.}
\prop{6.5}{t}{Ingen form er endeleg.}
\prop{6.51}{t}{Av 2.131 og 4.1: sidan seleksjonstrykka endrar seg, vil kvar realisert balanse med tida verte suboptimal\footnote{Den statiske koden må refaktorerast fordi trykka – brukartal, tryggleikstruslar, maskinvare – rører seg under føtene på arkitekturen.}.}
\prop{6.511}{t}{Kvar realisert form er gyldig berre for dei vilkåra som rådde i det augeblinket ho vart aktualisert.}
\prop{6.5111}{t}{Dei mest robuste tradisjonane er ikkje dei rigide, men dei med høgast agensstratifisering\footnote{Kompetansestrukturen i framtidsretta design kviler på å distribuere agentar slik at oppdateringar kan skje asynkront.}.}
\prop{6.51111}{o}{Fleire uavhengige tilpassingsmekanismar på fleire skalaer gjev raskare respons når landskapet endrar seg.}
\prop{6.51112}{t}{Det einaste varige er den generative strukturen: formrommet, seleksjonstrykka, landskapet, agentane.}
\prop{6.511121}{t}{Formene i seg sjølve er flyktige spor; grammatikken er stabil.}
\prop{6.6}{o}{Dette rammeverket er ikkje ein ontologisk påstand om kva form er.}
\prop{6.61}{o}{Det er eit koordinatsystem for spesifikasjon.}
\prop{6.611}{o}{At strukturen er substrat-uavhengig er eit teikn på fruktbarheit, ikkje eit prov for sanning\footnote{Det faktum at isomorfe formrom lèt seg etablere for stoldesign og for nettverksarkitektur bekreftar at teorien grip sjølve mekanismen bak strukturendring, uavhengig av ontologien til byggeklossane.}.}
\prop{6.6111}{o}{Agenten som skildrar formverda er sjølv ein del av ho.}
\prop{6.61111}{o}{Det finst ingen arkimedisk punkt på utsida.}
\prop{6.61112}{o}{Ei formlære er aldri transcendent; ho er alltid lokal og situert.}
\prop{6.61113}{o}{Å gje form er: å definere kva måltilstand materien skal navigere mot; å velje kva krefter som skal utgjere landskapet\footnote{I all software engineering og tenestedesign er rolla til arkitekten ikkje å diktere utsjånad, men å setje grensevilkåra slik at ei kompetent form sjølv kan emergere.}.}
\prop{6.611131}{o}{Det er òg å utvide agentens kognitive lyskjegle slik at latente affordansar vert synlege.}
\prop{6.611132}{o}{Det er å akseptere at forma emergerer mellom agentane.}
\prop{6.611133}{o}{Det er å leggje til rette for vilkåra som dei kompetente delane opererer under.}
\prop{6.7}{o}{Formverda er alt som er tilfelle. Tilpassingslandskapet er alt som verkar. Navigasjonen er utan slutt.}
\prop{6.71}{o}{Denne traktaten er sjølv ein posisjon i eit formrom.}
\prop{6.711}{o}{At teksten kan falsifiserast, er garantien for hennar gyldigheit\footnote{Også den vitskaplege metoden er her ein formalisme i eit epistemisk landskap; falsifiserbarheita tryggar at navigasjonen framleis er gyldig.}.}

